\documentclass[preprint,12pt]{elsarticle}

\usepackage{booktabs}
\usepackage[hidelinks]{hyperref}

\journal{Neuroscience Informatics}

\begin{document}
\begin{frontmatter}

\title{Duplicate-Aware Validation of Machine-Learning Models for Recorded Sleep-Status Labels in a Public Lifestyle Dataset}

\author[bracu]{Rifat Bin Reza\corref{cor1}}
\ead{rifat.bin.reza1@g.bracu.ac.bd}
\author[aau]{Md Ahadul Haque Suvo}
\author[bracu]{Atair Rahman Alvi}
\cortext[cor1]{Corresponding author}
\address[bracu]{BSRM School of Engineering, Department of Electrical and Electronic Engineering, BRAC University, Dhaka, Bangladesh}
\address[aau]{Department of Avionics Engineering, Aviation and Aerospace University, Dhaka, Bangladesh}

\begin{abstract}
Public tabular health datasets can give overly optimistic estimates when repeated predictor profiles cross validation boundaries. We re-evaluated Random Forest--XGBoost stacking for the Sleep Health and Lifestyle Dataset (374 rows) after auditing outcome coding and duplicate structure. The source file contains the literal outcome labels ``None'' (219), ``Insomnia'' (77), and ``Sleep Apnea'' (78), with no blank outcome values. After removing the administrative identifier and applying deterministic feature cleaning, 265 rows duplicated a retained predictor profile; 22 of 109 unique predictor profiles had conflicting recorded outcomes. We therefore used two repetitions of five-fold stratified group validation, keeping every exact predictor profile in one validation group. Within every outer-training partition, Random Forest and XGBoost probabilities for the logistic meta-learner were generated by four-fold group-aware out-of-fold fitting with preprocessing refit in each fold. No test-set tuning was performed. Mean outer-fold accuracy (standard deviation) was 0.887 (0.034) for stacking, 0.890 (0.033) for soft voting, and 0.896 (0.032) for logistic regression; corresponding macro-F1 values were 0.832, 0.835, and 0.852. The paired stacking-minus-soft-voting accuracy difference was -0.0031 (bootstrap 95\% interval -0.0162 to 0.0070); the first-repeat out-of-fold McNemar test was not significant ($p=0.6875$). These are internal estimates for dataset-specific recorded labels, not evidence of clinical screening or diagnosis.
\end{abstract}

\begin{keyword}
machine learning \sep tabular data \sep grouped validation \sep stacking \sep sleep health \sep reproducibility
\end{keyword}
\end{frontmatter}

\section{Introduction}

Small public tabular datasets are useful for teaching and benchmarking, but they are an unreliable setting for clinical claims. In particular, if identical or near-identical predictor records occur on both sides of a random train--test split, a model can appear to generalise while effectively seeing the same profile during training. This issue is especially consequential when the data-collection process and outcome construction are not documented.

Random Forest and gradient boosting are established methods for heterogeneous tabular data \cite{breiman2001,chen2016}, and stacking learns a second-level combination of base-model predictions \cite{wolpert1992}. The methodological contribution of this revision is not a new classifier. It is a fully executable audit and evaluation protocol for this specific public benchmark: literal outcome verification, duplicate-aware grouping, fold-specific preprocessing, out-of-fold meta-features, matched comparators, and release of all predictions and split assignments. The analysis asks a deliberately narrow question: how do fixed, conventional models discriminate the dataset's \emph{recorded} labels when exact predictor duplicates cannot cross a validation boundary?

\section{Materials and methods}

\subsection{Data audit and target definition}

The analysis used the CSV distributed with the Sleep Health and Lifestyle Dataset \cite{kaggle_sleep}. The file was read with default missing-value conversion disabled. The outcome column was stripped of surrounding whitespace and required to contain only the three literal source strings shown in Table~\ref{tab:audit}; blank, missing, or undocumented values would stop the analysis. Thus, ``None'' is not an imputation or a conversion of missing data. Its clinical meaning remains unknown: it denotes only the source-file category and cannot be treated as independently adjudicated absence of disease.

\begin{table}[t]
\centering
\caption{Source-data and duplicate-structure audit after deterministic feature preparation.}
\label{tab:audit}
\begin{tabular}{lr}
\toprule
Item & Count \\
\midrule
Rows & 374 \\
Literal ``None'' outcomes & 219 \\
Literal ``Insomnia'' outcomes & 77 \\
Literal ``Sleep Apnea'' outcomes & 78 \\
Blank or missing outcomes & 0 \\
Unique retained predictor profiles & 109 \\
Rows duplicating a retained predictor profile & 265 \\
Predictor profiles with conflicting recorded outcomes & 22 \\
\bottomrule
\end{tabular}
\end{table}

Person ID was removed before modelling. ``Normal Weight'' was harmonised to ``Normal'', and Blood Pressure was parsed into numerical systolic and diastolic components. The retained profile was defined by all remaining predictor values after these deterministic steps. Group identifiers were created by factorising the complete predictor tuples, not by outcome. Consequently, profiles with conflicting labels were still held together, preventing the same predictors from appearing in both development and evaluation sets.

\subsection{Models and validation}

Numerical variables were standardised and categorical variables one-hot encoded. Each transformer was fit only on the relevant training data. The primary stack used fixed specifications: a Random Forest (150 trees; minimum leaf size 2), XGBoost (150 trees; learning rate 0.03; maximum depth 3; subsample and column subsample 0.9), and a multinomial logistic meta-learner ($C=1.0$). The meta-learner received only base-model probability vectors; raw-feature pass-through was not used.

For each of two fixed shuffled repetitions (seeds 17 and 29), StratifiedGroupKFold produced five outer folds. Every exact predictor-profile group was confined to one outer test fold. Inside each outer training partition, four-fold StratifiedGroupKFold generated out-of-fold base probabilities. Each base pipeline, including its encoder and scaler, was freshly fit within that inner-fold training partition. The meta-learner was fit on these out-of-fold probabilities; final base learners were then refit on the complete outer training data before predicting the outer test fold. Hyperparameters were fixed before this evaluation; no test-fold results were used to tune a model.

The matched comparators were soft voting of the same two base learners, each base learner, logistic regression, a shallow decision tree, and histogram gradient boosting. We report accuracy, macro-F1, multiclass log loss, and multiclass Brier score for each of ten outer folds. The released prediction file includes the reference label, predicted label, all class probabilities, group ID, fold, and repetition for every model. The two repetitions share records and are correlated; their variation is an internal robustness description, not an independent confidence interval or external validation.

\subsection{Statistical comparison and calibration}

For the prespecified stacking versus soft-voting comparison, we calculated the mean matched outer-fold accuracy difference and a 10,000-resample percentile bootstrap interval over the ten fold-level differences. We also applied an exact McNemar test to the first repetition's complete out-of-fold predictions. This test is descriptive because the data remain one public dataset. Multiclass Brier scores and class-wise reliability curves from the first repetition's out-of-fold probabilities are supplied as reproducibility artefacts. Post-hoc SHAP, LIME, PCA, and ALE claims from the earlier exploratory notebooks were excluded: the duplicate structure, conflicting labels, and correlated predictors make unvalidated feature rankings too unstable for a substantive interpretation in this revision.

\section{Results}

Table~\ref{tab:performance} reports the ten grouped outer-fold estimates. Logistic regression had the highest mean accuracy and macro-F1. Stacking had a slightly lower mean accuracy than soft voting, although its mean log loss and Brier score were modestly lower. The observed difference does not establish that either ensemble is superior.

\begin{table}[t]
\centering
\caption{Repeated grouped outer-validation performance. Values are mean (standard deviation) across ten outer folds. Lower is better for log loss and multiclass Brier score.}
\label{tab:performance}
\resizebox{\textwidth}{!}{%
\begin{tabular}{lrrrr}
\toprule
Model & Accuracy & Macro-F1 & Log loss & Brier \\
\midrule
Logistic regression & 0.896 (0.032) & 0.852 (0.073) & 0.441 (0.117) & 0.190 (0.048) \\
Shallow decision tree & 0.858 (0.037) & 0.807 (0.063) & 0.933 (0.635) & 0.227 (0.057) \\
Histogram gradient boosting & 0.861 (0.067) & 0.800 (0.083) & 0.620 (0.211) & 0.239 (0.094) \\
Random Forest & 0.886 (0.034) & 0.837 (0.060) & 0.679 (0.331) & 0.192 (0.049) \\
XGBoost & 0.889 (0.032) & 0.832 (0.070) & 0.409 (0.108) & 0.191 (0.057) \\
Soft voting (RF + XGBoost) & 0.890 (0.033) & 0.835 (0.078) & 0.422 (0.105) & 0.189 (0.051) \\
Stacking (RF + XGBoost) & 0.887 (0.034) & 0.832 (0.069) & 0.397 (0.094) & 0.187 (0.054) \\
\bottomrule
\end{tabular}
}
\end{table}

The stacking-minus-soft-voting mean accuracy difference was -0.0031. Its bootstrap 95\% interval was -0.0162 to 0.0070 and therefore crossed zero. In the first repetition, stacking alone was correct for two observations and soft voting alone for four; the exact McNemar $p$ value was 0.6875. The small number of discordant predictions and the limited sample make this comparison low-powered. It supports neither a stacking advantage nor an equivalence claim.

\section{Discussion}

The change from the earlier single random hold-out result illustrates why validation design matters more than marginal differences among conventional classifiers on this dataset. After duplicate-profile protection, all models produced lower and more variable estimates, and a simple logistic regression was competitive with or better than the ensemble on discrimination metrics. The numerical advantage of stacking on probability scores is small and must not be interpreted as evidence of clinical usefulness.

The target is a source-file label, not an independently confirmed diagnosis. In addition, sleep duration and sleep quality may overlap conceptually with symptoms used to identify insomnia, and 22 identical predictor profiles carry different labels. A high internal score could therefore reflect dataset construction, label conventions, or predictor--outcome overlap rather than transferable physiological signal. A computational result on these labels is not a real-world screening result: clinical screening would require a prospectively specified intended use, independently adjudicated reference standards, calibration in the target population, subgroup assessment, and external validation in new clinical cohorts.

\section{Limitations and conclusion}

This study remains limited by one small public dataset, undocumented data provenance, non-adjudicated outcome meaning, repeated folds that are not independent datasets, and the absence of external clinical validation. We deliberately removed unsupported local and global explanation claims rather than presenting unstable attributions as mechanisms. Future work should use prospectively collected cohorts, prespecified group definitions, nested hyperparameter tuning confined to development data, calibration assessment with adequate sample sizes, and external validation.

Under duplicate-aware validation, fixed RF--XGBoost stacking did not outperform its soft-voting comparator or logistic regression on mean accuracy or macro-F1. The contribution is a transparent, reproducible benchmark audit, not a proposed diagnostic system.

\section*{Declaration of competing interest}
The authors declare no competing interests.

\section*{Data and code availability}
The public dataset is cited above. The accompanying repository contains \texttt{revised\_analysis.py}, package requirements, and regenerated files in \texttt{revision\_results/}: software versions, exact split assignments, fold metrics, observation-level predictions and probabilities, paired-comparison results, and calibration artefacts. Run the analysis with \texttt{python} followed by \texttt{revised\_analysis.py}.

\bibliographystyle{elsarticle-num}
\bibliography{revised_references}
\end{document}
